% ============================================================
% DOCUMENTCLASS
% ============================================================
\documentclass[12pt,oneside]{book}

% ============================================================
% METADATOS
% ============================================================
\newcommand{\universidad}{Universidad de Buenos Aires (UBA)}
\newcommand{\facultad}{Facultad de Derecho}
\newcommand{\materia}{Concursos y Quiebras}
\newcommand{\titulo}{Efectos de la quiebra, continuación de la explotación de la empresa y extensión de quiebra}
\newcommand{\autor}{Isaac Edgar Camacho Ocampo}
\newcommand{\anio}{2026}

% ============================================================
% PAQUETES
% ============================================================
\usepackage[utf8]{inputenc}
\usepackage[T1]{fontenc}
\usepackage[spanish,es-nodecimaldot]{babel}

\usepackage[
    a4paper,
    top=2.5cm,
    bottom=2.5cm,
    left=3cm,
    right=2.5cm,
    headheight=15pt
]{geometry}

\usepackage{amsmath}
\usepackage{amssymb}
\usepackage{mathtools}

\usepackage{graphicx}
\usepackage{float}
\usepackage{caption}
\usepackage{subcaption}

\usepackage{booktabs}
\usepackage{array}
\usepackage{longtable}
\usepackage{tabularx}

\usepackage{enumitem}

\usepackage{xcolor}
\usepackage[most]{tcolorbox}

\usepackage{fancyhdr}
\usepackage{ragged2e}

\usepackage[
    colorlinks=true,
    linkcolor=blue!50!black,
    citecolor=green!40!black,
    urlcolor=blue!60!black,
    pdftitle={\titulo},
    pdfauthor={\autor},
    pdfsubject={Apuntes universitarios},
    bookmarks=true
]{hyperref}

% ============================================================
% CONFIGURACIÓN
% ============================================================
\graphicspath{
    {./img/}
    {./graficos/}
    {./figuras/}
}

\setcounter{tocdepth}{2}

\setlength{\parindent}{0pt}
\setlength{\parskip}{0.5em}

\setlist[itemize]{
    topsep=0.3em,
    itemsep=0.2em
}

\setlist[enumerate]{
    topsep=0.3em,
    itemsep=0.2em
}

\clubpenalty=10000
\widowpenalty=10000

\pagestyle{fancy}
\fancyhf{}
\fancyhead[L]{\nouppercase{\leftmark}}
\fancyhead[R]{\materia}
\fancyfoot[C]{\thepage}
\renewcommand{\headrulewidth}{0.4pt}
\renewcommand{\footrulewidth}{0pt}

\fancypagestyle{plain}{
    \fancyhf{}
    \fancyfoot[C]{\thepage}
    \renewcommand{\headrulewidth}{0pt}
}

% ============================================================
% COMANDOS
% ============================================================
\newcommand{\abs}[1]{\left|#1\right|}
\newcommand{\norm}[1]{\left\lVert#1\right\rVert}
\newcommand{\vect}[1]{\mathbf{#1}}
\newcommand{\deriv}[2]{\frac{d#1}{d#2}}
\newcommand{\pderiv}[2]{\frac{\partial#1}{\partial#2}}

% Tipos de columna y anchos auxiliares
\newcolumntype{Y}{>{\raggedright\arraybackslash}X}
\newcommand{\cornellfull}{\dimexpr0.91\textwidth+2\tabcolsep\relax}

% ============================================================
% ENTORNOS / CAJAS
% ============================================================
\newtcolorbox{definition}[1]{
    enhanced,
    breakable,
    title={Definición: #1},
    fonttitle=\bfseries,
    colback=blue!3,
    colframe=blue!50!black,
    boxrule=0.8pt,
    arc=2mm
}

\newtcolorbox{important}{
    enhanced,
    breakable,
    title={Importante},
    fonttitle=\bfseries,
    colback=yellow!8,
    colframe=orange!70!black,
    boxrule=0.8pt,
    arc=2mm
}

\newtcolorbox{example}[1]{
    enhanced,
    breakable,
    title={Ejemplo: #1},
    fonttitle=\bfseries,
    colback=green!4,
    colframe=green!40!black,
    boxrule=0.8pt,
    arc=2mm
}

\newtcolorbox{formula}[1]{
    enhanced,
    breakable,
    title={Fórmula / Regla: #1},
    fonttitle=\bfseries,
    colback=purple!4,
    colframe=purple!50!black,
    boxrule=0.8pt,
    arc=2mm
}

\newtcolorbox{exercise}[1]{
    enhanced,
    breakable,
    title={Ejercicio: #1},
    fonttitle=\bfseries,
    colback=gray!5,
    colframe=gray!60!black,
    boxrule=0.8pt,
    arc=2mm
}

\newtcolorbox{note}{
    enhanced,
    breakable,
    title={Nota},
    fonttitle=\bfseries,
    colback=white,
    colframe=gray!50,
    boxrule=0.6pt,
    arc=2mm
}

% Caja para normas legales citadas en el apunte (regla de derecho)
\newtcolorbox{norma}[1]{
    enhanced,
    breakable,
    title={Norma: #1},
    fonttitle=\bfseries,
    fontupper=\small,
    colback=yellow!4,
    colframe=yellow!45!black,
    boxrule=0.8pt,
    arc=2mm
}

% ============================================================
% CONFIGURACIÓN FINAL
% ============================================================
\captionsetup{font=small,labelfont=bf}
\renewcommand{\arraystretch}{1.25}

% ============================================================
% DOCUMENT
% ============================================================
\begin{document}

% ============================================================
% PORTADA
% ============================================================
\begin{titlepage}

\thispagestyle{empty}

\begin{center}

\vspace*{1cm}

{\large \textsc{\universidad}}

\vspace{0.2cm}

{\large \textsc{\facultad}}

\vspace{3cm}

{\Huge \textbf{\materia}}

\vspace{1cm}

{\LARGE \titulo}

\vspace{0.8cm}

\textit{Comprender cada concepto es el primer paso para convertir el estudio en conocimiento.}

\vspace{1.2cm}

\rule{0.70\textwidth}{0.5pt}

\vspace{0.5cm}

{\large Apuntes de estudio}

\vspace{0.5cm}

\rule{0.70\textwidth}{0.5pt}

\vfill

{\large \autor}

\vspace{0.3cm}

{\large \anio}

\end{center}

\vfill

\begin{flushleft}
\textit{Este es un modesto aporte a los estudiantes de ``\materia'' no pretende sustituir el material oficial de la materia.}
\end{flushleft}

\end{titlepage}

% ============================================================
% ÍNDICE
% ============================================================
\frontmatter
\tableofcontents
\mainmatter

% ============================================================
% CONTENIDO
% ============================================================

% ------------------------------------------------------------
\chapter*{Visión general}
\markboth{Visión general}{Visión general}
\addcontentsline{toc}{chapter}{Visión general}
% ------------------------------------------------------------

Este material estudia qué ocurre desde el momento en que se decreta la quiebra. En primer lugar se repasa cómo se llega a ella (de manera directa, indirecta o por extensión) y cómo se la ataca. Luego se recorren sus efectos: sobre la \textbf{persona del fallido} (inhabilitación, deber de cooperación, viajes), sobre su \textbf{patrimonio} (desapoderamiento), sobre los \textbf{actos realizados en el período de sospecha} (arts.~118 a 120), sobre los \textbf{contratos vigentes} (arts.~143 a 145 y 157) y sobre los \textbf{trabajadores}. A continuación se estudia la \textbf{continuación de la explotación de la empresa} (arts.~189 y siguientes) y, por último, la \textbf{extensión de quiebra} (arts.~160 y 161), junto con la forma en que los acreedores verifican sus créditos.

\begin{example}{la panadería del barrio (ejemplo elaborado para unificar los bloques)}
Una panadería de barrio deja de poder pagar a sus proveedores. Los bloques de contenido pueden leerse como las preguntas que se formulan, en orden, cuando esa panadería quiebra:
\begin{enumerate}
    \item \textbf{Repaso: ¿cómo llegó a la quiebra?} El dueño la pidió, un proveedor la pidió o fracasó un concurso previo.
    \item \textbf{Persona y patrimonio: ¿qué le pasa al dueño y a sus bienes?} Queda inhabilitado, colabora con el síndico y el síndico administra el local y el horno.
    \item \textbf{Actos perjudiciales: ¿qué hizo el dueño antes de cerrar?} Si le vendió el local a su socio o pagó por adelantado a un amigo, esos actos pueden caerse.
    \item \textbf{Relaciones preexistentes: ¿qué pasa con el alquiler del local, los contratos con proveedores y los empleados?}
    \item \textbf{Continuación: ¿puede la panadería seguir horneando si es útil para el barrio y no genera más deudas?} Quizás los empleados la sigan como cooperativa.
    \item \textbf{Extensión: ¿y si el dueño mezclaba sus cuentas con las de otra empresa?} Entonces la quiebra puede alcanzar también a esa otra.
\end{enumerate}
Hilo conductor: cada bloque responde a una pregunta práctica sobre la misma panadería, desde cómo quiebra hasta quién más puede quedar alcanzado.
\end{example}

% ------------------------------------------------------------
\chapter[La quiebra y su origen]{Repaso: la quiebra y su origen}
% ------------------------------------------------------------

\section{La quiebra como proceso de ejecución colectiva}

La quiebra es un \textbf{proceso de ejecución colectiva}. Requiere la concurrencia de los dos presupuestos, no exige pluralidad de acreedores y presenta los mismos caracteres que todos los procesos concursales: \textbf{universalidad}, \textbf{igualdad} y \textbf{colectividad}.

\section{Origen de la quiebra}

\subsection{Clasificación}

Por su origen se distinguen la \textbf{quiebra directa}, la \textbf{quiebra indirecta} y la \textbf{extensión de quiebra}. La quiebra directa puede ser pedida por el deudor o por el acreedor. La quiebra indirecta se produce cuando fracasa el concurso preventivo.

\subsection{Quiebra indirecta}

La quiebra indirecta se da por la frustración del concurso preventivo, y las causales se encuentran mencionadas en la propia ley: que el deudor no haya presentado acuerdo, que no haya logrado las conformidades o que el salvataje no haya tenido un buen resultado.

\subsection{Quiebra pedida por el acreedor}

La \textbf{quiebra directa necesaria} se da cuando la pide el acreedor, quien debe cumplir ciertos requisitos: demostrar que su crédito es exigible, que el deudor es un sujeto concursable y que se encuentra en cesación de pagos. El juez puede ordenar medidas sumarias, luego se cita al deudor y finalmente el juez resuelve si decreta o no la quiebra.

\subsection{Quiebra pedida por el deudor}

El deudor puede pedir su propia quiebra (\textbf{quiebra directa voluntaria}). Debe cumplir con todos los requisitos del artículo 11 y, si le falta alguno, ello no obsta a que se le decrete la quiebra. Tanto el acreedor como el deudor pueden desistir del pedido de quiebra en los momentos que la ley prevé.

\begin{norma}{Art.~11 · Requisitos de la petición del deudor}
Citado respecto del pedido de quiebra por el propio deudor: debe cumplir los requisitos del artículo 11 y, si falta alguno, ello no obsta a que se decrete la quiebra. Entre ellos, el inciso 8, incorporado por la Ley 26.684, dispone:

«8) Acompañar nómina de empleados, con detalle de domicilio, categoría, antigüedad y última remuneración recibida. Deberá acompañarse también declaración sobre la existencia de deuda laboral y de deuda con los organismos de la seguridad social certificada por contador público.»

\par\smallskip\textit{Inciso 8. Texto literal verificado. Ley 24.522 según Ley 26.684 (B.O. 30/06/2011).}
\end{norma}

\begin{table}[H]
\centering
\small
\caption{Vías de origen de la quiebra}
\begin{tabularx}{\textwidth}{>{\raggedright\arraybackslash}p{0.2\textwidth} >{\raggedright\arraybackslash}p{0.25\textwidth} Y}
\toprule
\textbf{Vía} & \textbf{Quién la origina} & \textbf{Rasgos señalados} \\
\midrule
\textbf{Directa voluntaria} & El deudor (pide su propia quiebra) & Debe cumplir los requisitos del art.~11; si falta alguno, igual puede decretarse la quiebra. \\
\textbf{Directa necesaria} & El acreedor & Debe demostrar crédito exigible, deudor sujeto concursable y cesación de pagos; el juez puede ordenar medidas sumarias y se cita al deudor. \\
\textbf{Indirecta} & Fracaso del concurso preventivo & No se presentó acuerdo, no se lograron las conformidades o el salvataje no tuvo buen resultado. \\
\textbf{Extensión de quiebra} & Quiebra de otro sujeto (quiebra principal) & Se estudia en el último capítulo (arts.~160 y 161). \\
\bottomrule
\end{tabularx}
\end{table}

\section{Sentencia de quiebra y recursos}

\subsection{Contenido de la sentencia}

El contenido de la sentencia que decreta la quiebra está fijado en el artículo 88.

\begin{norma}{Art.~88 · Contenido de la sentencia de quiebra}
Norma citada como la que fija los contenidos de la sentencia que decreta la quiebra.

\par\smallskip\textit{El texto literal debe cotejarse con la versión vigente de la Ley 24.522; aquí solo se registra lo expuesto.}
% TODO: contenido incompleto o ambiguo en las notas originales (art. 88 solo mencionado)
\end{norma}

\subsection{Recursos contra el decreto de quiebra}

Los recursos contra el decreto de quiebra conforman un \textbf{sistema recursivo propio} de la Ley de Concursos y Quiebras. Aunque el Código Procesal también regula recursos, los que proceden contra el decreto de quiebra son los establecidos en la propia ley concursal: la \textbf{conversión}, la \textbf{reposición} y la \textbf{incompetencia}.

\begin{itemize}
    \item \textbf{Conversión:} permite transformar la quiebra en un concurso preventivo. Hay casos en que procede y casos en que no, y se exigen requisitos.
    % TODO: contenido incompleto o ambiguo en las notas originales (requisitos de la conversión no desarrollados)
    \item \textbf{Reposición:} se funda en que no se dan algunos de los presupuestos, ya sea porque el deudor no está en cesación de pagos o porque no es sujeto concursable. Existen dos tipos: \textbf{con trámite} y \textbf{sin trámite}. La reposición sin trámite se da cuando el deudor deposita para desvirtuar que está en cesación de pagos.
    \item \textbf{Incompetencia:} no deja sin efecto el decreto de quiebra; el juez declara su incompetencia y la actuación se remite al juez competente.
\end{itemize}

\begin{table}[H]
\centering
\small
\caption{Recursos contra el decreto de quiebra}
\begin{tabularx}{\textwidth}{>{\raggedright\arraybackslash}p{0.22\textwidth} >{\raggedright\arraybackslash}p{0.30\textwidth} Y}
\toprule
\textbf{Recurso} & \textbf{Qué busca} & \textbf{Precisión} \\
\midrule
\textbf{Conversión} & Transformar la quiebra en concurso preventivo & Hay casos en que procede y casos en que no; se exigen requisitos. \\
\textbf{Reposición con trámite} & Atacar la quiebra por falta de presupuestos & Se funda en que no hay cesación de pagos o el deudor no es sujeto concursable. \\
\textbf{Reposición sin trámite} & Desvirtuar la cesación de pagos & El deudor deposita para desvirtuar que está en cesación. \\
\textbf{Incompetencia} & Remitir la causa al juez competente & No deja sin efecto el decreto de quiebra. \\
\bottomrule
\end{tabularx}
\end{table}

\subsection*{Cuadro Cornell del capítulo}

\begin{longtable}{>{\raggedright\arraybackslash}p{0.27\textwidth} >{\raggedright\arraybackslash}p{0.64\textwidth}}
\toprule
\textbf{Preguntas / palabras clave} & \textbf{Notas / respuestas} \\
\midrule
\endfirsthead
\toprule
\textbf{Preguntas / palabras clave} & \textbf{Notas / respuestas} \\
\midrule
\endhead
\textbf{¿Qué caracteriza a la quiebra como proceso?} & Es un proceso de ejecución colectiva que requiere los dos presupuestos, no exige pluralidad de acreedores y presenta los caracteres de todos los procesos concursales: universalidad, igualdad y colectividad. \\[0.6em]
\textbf{¿Por qué vías se llega a la quiebra?} & Quiebra directa (la pide el deudor o el acreedor), quiebra indirecta (fracasa el concurso preventivo) y extensión de quiebra. \\[0.6em]
\textbf{¿Qué debe acreditar el acreedor que pide la quiebra?} & Que su crédito es exigible, que el deudor es sujeto concursable y que está en cesación de pagos. El juez puede ordenar medidas sumarias, se cita al deudor y luego resuelve. \\[0.6em]
\textbf{¿Qué requisitos debe cumplir el deudor que pide su propia quiebra?} & Los del art.~11 (entre ellos, el inciso 8 sobre nómina de empleados y declaración de deuda laboral y previsional certificada por contador público). Si falta alguno, igual puede decretarse la quiebra. \\[0.6em]
\textbf{¿Qué recursos existen contra el decreto de quiebra?} & Conversión (pasar a concurso preventivo), reposición con trámite (faltan presupuestos), reposición sin trámite (el deudor deposita para desvirtuar la cesación de pagos) e incompetencia (remite al juez competente sin dejar sin efecto la quiebra). \\
\midrule
\multicolumn{2}{>{\raggedright\arraybackslash}p{\cornellfull}}{\textbf{Resumen:} la quiebra es un proceso de ejecución colectiva con tres caracteres (universalidad, igualdad, colectividad). Se origina de forma directa (deudor o acreedor), indirecta (fracaso del concurso preventivo) o por extensión, y se la ataca mediante los recursos propios de la ley concursal: conversión, reposición e incompetencia.} \\
\bottomrule
\end{longtable}

% ------------------------------------------------------------
\chapter[Efectos sobre persona y patrimonio]{Efectos de la quiebra sobre la persona y el patrimonio}
% ------------------------------------------------------------

\section{Mapa de los efectos de la quiebra}

Los efectos de la quiebra son aquellos que se producen una vez que la quiebra ya fue decretada. Se los clasifica bajo distintos ejes:

\begin{table}[H]
\centering
\small
\caption{Ejes de los efectos de la quiebra}
\begin{tabularx}{\textwidth}{>{\raggedright\arraybackslash}p{0.32\textwidth} Y}
\toprule
\textbf{Eje} & \textbf{Contenido} \\
\midrule
\textbf{1. Persona del fallido} & Inhabilitación, deber de cooperación, viajes, empleo, profesión u oficio \\
\textbf{2. Patrimonio} & Desapoderamiento, clausura, inventario, bienes exceptuados, herencias y donaciones \\
\textbf{3. Actos perjudiciales a los acreedores} & Período de sospecha, ineficacia de pleno derecho (art.~118), ineficacia por conocimiento (art.~119), acción (art.~120) \\
\textbf{4. Relaciones jurídicas preexistentes} & Fuero de atracción, intereses, restitución, contratos en curso, locación \\
\textbf{5. Contratos de trabajo} & Suspensión por 60 días; continuidad o extinción según la continuación de la empresa \\
\bottomrule
\end{tabularx}
\end{table}

\section{Efectos sobre la persona del fallido}

\subsection{Inhabilitación}

La \textbf{inhabilitación} es un efecto muy importante de la quiebra. Se advierte que no se encuentra regulada dentro de la parte de la ley dedicada a los efectos de la quiebra, sino al final.
% TODO: contenido incompleto o ambiguo en las notas originales (ubicación normativa de la inhabilitación)

Cuando la persona que quiebra es una \textbf{persona jurídica}, la sociedad se disuelve y se liquida, y la inhabilitación es \textbf{definitiva}. En cambio, la \textbf{persona humana} queda inhabilitada por el término de \textbf{un año}. Esa inhabilitación implica que no puede ejercer el comercio ni realizar algunas funciones dentro de sociedades. Se trata de un plazo corto, que puede acortarse o ampliarse según exista o no causa penal.

\begin{table}[H]
\centering
\small
\caption{Inhabilitación del fallido}
\begin{tabularx}{\textwidth}{>{\raggedright\arraybackslash}p{0.22\textwidth} >{\raggedright\arraybackslash}p{0.34\textwidth} Y}
\toprule
\textbf{Sujeto} & \textbf{Plazo} & \textbf{Alcance} \\
\midrule
\textbf{Persona jurídica} & Definitiva (la sociedad se disuelve y se liquida) & --- \\
\textbf{Persona humana} & Un año (puede acortarse o ampliarse según haya causa penal) & No puede ejercer el comercio ni realizar algunas funciones dentro de sociedades. \\
\bottomrule
\end{tabularx}
\end{table}

\subsection{Deber de colaboración o cooperación y viajes}

La quiebra es un proceso de liquidación. En la medida en que el juez y el síndico necesitan liquidar los bienes, requieren la colaboración del propio deudor para que ponga los bienes a disposición e indique dónde se encuentran. Por ello se le exige \textbf{autorización del juzgado para viajar al exterior}.

En el concurso preventivo, si el viaje es por menos de 40 días el deudor debe informar al juez, y por más de 40 días necesita autorización. En la quiebra la restricción es más intensa: el fallido no puede salir del país hasta que se haya presentado el \textbf{informe general del síndico}, y el juez puede ampliar ese plazo de interdicción por \textbf{seis meses más} después del informe, si corresponde. También puede pedirse autorización para casos particulares. La interdicción de salida del país se relaciona con el deber de colaboración o cooperación.
% TODO: contenido incompleto o ambiguo en las notas originales (el apunte señala que la ley no fue consultada durante la exposición)

\subsection{Empleo, profesión u oficio}

Si bien el fallido no puede ejercer el comercio, sí puede desempeñarse en su empleo, su profesión o su oficio: la ley le permite realizar trabajos, pues el fallido necesita medios para vivir.

\section{Desapoderamiento y bienes del fallido}

\subsection{Imposibilidad de disponer, clausura e inventario}

El fallido no puede disponer de sus bienes: lo hace el \textbf{síndico}, dado que el objetivo de la quiebra es liquidar los bienes del deudor. Al decretarse la quiebra, lo primero que se hace es la \textbf{clausura del establecimiento}. Luego se realiza un \textbf{inventario} de bienes y de documentación, y se incautan esos bienes y papeles, que quedan en poder del síndico para proceder después a su liquidación.

\begin{definition}{Desapoderamiento}
Efecto de la quiebra por el cual el fallido pierde la administración y la disposición de sus bienes, que pasan a ser administrados por el síndico. Se distingue del concurso preventivo, donde la administración de los bienes sigue en cabeza del propio deudor.
\end{definition}

\begin{table}[H]
\centering
\small
\caption{Administración de los bienes}
\begin{tabularx}{\textwidth}{>{\raggedright\arraybackslash}p{0.24\textwidth} >{\raggedright\arraybackslash}p{0.30\textwidth} Y}
\toprule
\textbf{Proceso} & \textbf{Quién administra} & \textbf{Precisión} \\
\midrule
\textbf{Concurso preventivo} & El propio deudor & Conserva la administración de sus bienes. \\
\textbf{Quiebra} & El síndico (desapoderamiento) & El fallido no dispone de sus bienes; clausura, inventario e incautación de bienes y documentación para liquidar. \\
\bottomrule
\end{tabularx}
\end{table}

\subsection{Bienes exceptuados}

Existen bienes exceptuados por la propia ley. Por ejemplo, quedan exceptuados los bienes de \textbf{protección de la familia} y los bienes que sirven al deudor para ejercer su profesión u oficio.

\subsection{Herencias y donaciones}

El fallido difícilmente querrá iniciar una sucesión para cobrar una herencia, porque esos bienes ingresarían al patrimonio que se terminará liquidando. Por eso, muchas veces los fallidos terminan rechazando las herencias. Sin embargo, interviene el síndico, pues lo que se pretende es que ingresen todos esos bienes.

\subsection*{Cuadro Cornell del capítulo}

\begin{longtable}{>{\raggedright\arraybackslash}p{0.27\textwidth} >{\raggedright\arraybackslash}p{0.64\textwidth}}
\toprule
\textbf{Preguntas / palabras clave} & \textbf{Notas / respuestas} \\
\midrule
\endfirsthead
\toprule
\textbf{Preguntas / palabras clave} & \textbf{Notas / respuestas} \\
\midrule
\endhead
\textbf{¿Qué ejes organizan los efectos de la quiebra?} & Persona del fallido, patrimonio, actos perjudiciales a los acreedores, relaciones jurídicas preexistentes y contratos de trabajo. \\[0.6em]
\textbf{¿Qué plazo de inhabilitación rige?} & Definitiva para las personas jurídicas (la sociedad se disuelve y se liquida); un año para las personas humanas, plazo que puede acortarse o ampliarse según haya causa penal. No puede ejercer el comercio ni algunas funciones dentro de sociedades. \\[0.6em]
\textbf{¿Qué exige el deber de cooperación?} & Que el fallido ponga los bienes a disposición e indique dónde se encuentran. Se relaciona con la necesidad de autorización para viajar y con la interdicción de salida del país hasta el informe general del síndico (ampliable por seis meses más). \\[0.6em]
\textbf{¿Puede el fallido trabajar?} & Sí: aunque no puede ejercer el comercio, puede desempeñarse en su empleo, profesión u oficio. \\[0.6em]
\textbf{¿Qué es el desapoderamiento?} & Pérdida de la administración y disposición de los bienes, que pasan al síndico. Se clausura el establecimiento, se inventaría e incauta. En el concurso preventivo, en cambio, la administración sigue en cabeza del deudor. \\[0.6em]
\textbf{¿Qué bienes quedan exceptuados y qué ocurre con herencias y donaciones?} & Se exceptúan los bienes de protección de la familia y los que sirven para ejercer la profesión u oficio. Las herencias y donaciones ingresan al patrimonio a liquidar, con intervención del síndico. \\
\midrule
\multicolumn{2}{>{\raggedright\arraybackslash}p{\cornellfull}}{\textbf{Resumen:} decretada la quiebra, el fallido queda inhabilitado, debe cooperar con el juez y el síndico, y pierde la administración y disposición de sus bienes (desapoderamiento). Conserva la posibilidad de trabajar en su profesión u oficio y se exceptúan ciertos bienes; todo lo demás, incluidas herencias y donaciones, se destina a la liquidación.} \\
\bottomrule
\end{longtable}

% ------------------------------------------------------------
\chapter[Actos perjudiciales a los acreedores]{Actos perjudiciales a los acreedores}
% ------------------------------------------------------------

\section{Período de sospecha}

\begin{definition}{Período de sospecha}
Período de tiempo determinado por la ley para establecer si hay actos realizados por el deudor que puedan considerarse ineficaces. Comprende desde el inicio de la cesación de pagos hasta el decreto de quiebra.
\end{definition}

Es uno de los aspectos más importantes de los efectos de la quiebra. Todo el lapso que va desde el inicio de la cesación de pagos hasta el decreto de quiebra se considera período de sospecha. La línea de tiempo difiere según se trate de una quiebra directa o de una quiebra indirecta.
% TODO: contenido incompleto o ambiguo en las notas originales (el audio de este tramo es confuso; no se detalla cómo difiere la línea de tiempo)

\section{Actos ineficaces de pleno derecho (art.~118)}

Dentro del período de sospecha se detectan los \textbf{actos ineficaces}. Los dos artículos centrales son el \textbf{art.~118} (actos ineficaces de pleno derecho) y el \textbf{art.~119} (actos ineficaces por conocimiento del estado de cesación de pagos).

Los actos ineficaces de pleno derecho son aquellos en los que \textbf{no se necesita demostrar nada}: el propio juez, de oficio, al detectarlos, los considera ineficaces. Son los siguientes:

\begin{itemize}
    \item \textbf{Actos a título gratuito.}
    \item \textbf{Pagos realizados de forma anticipada}, antes de que venza el crédito.
    \item Constitución de una \textbf{garantía real} a un acreedor que era quirografario, que pasa así a ser privilegiado. Con ello se desafecta un bien que garantizaba a todos los acreedores para beneficiar a uno.
\end{itemize}

\begin{norma}{Art.~118 · Actos ineficaces de pleno derecho}
Contenido expuesto: son ineficaces respecto de los acreedores, sin necesidad de probar nada y declarables de oficio por el juez, los actos realizados por el deudor en el período de sospecha que consistan en (1) actos a título gratuito, (2) pagos anticipados de deudas aún no vencidas y (3) constitución de garantía real en favor de un acreedor quirografario, que pasa así a ser privilegiado.

\par\smallskip\textit{Texto literal pendiente de verificación: se resume lo expuesto; cotejar con la versión vigente de la Ley 24.522.}
% TODO: contenido incompleto o ambiguo en las notas originales (texto literal no verificado)
\end{norma}

\section{Actos ineficaces por conocimiento de la cesación de pagos (art.~119)}

El art.~119 se refiere a aquellos actos en los que el \textbf{tercero que interviene en la operación conocía que el deudor estaba en cesación de pagos} y que terminan afectando a todos los acreedores.

\begin{example}{venta del inmueble al socio}
El deudor tenía un inmueble y se lo vende a su socio. Si el producido de la venta desaparece (no reaparece bajo otra forma en el patrimonio del deudor), es como si el inmueble hubiera desaparecido, y los acreedores resultan afectados. En cambio, si con lo obtenido el deudor compra otro inmueble o deposita el dinero en una cuenta bancaria (por ejemplo, porque quería pagar a sus acreedores), no hay afectación de los acreedores y, por lo tanto, no puede hablarse de acto ineficaz.
\end{example}

\begin{norma}{Art.~119 · Actos ineficaces por conocimiento de la cesación de pagos}
Contenido expuesto: los demás actos perjudiciales para los acreedores, realizados en el período de sospecha, pueden ser declarados ineficaces si el tercero que contrató con el deudor conocía el estado de cesación de pagos.

\par\smallskip\textit{Texto literal pendiente de verificación: se resume lo expuesto; cotejar con la versión vigente de la Ley 24.522.}
% TODO: contenido incompleto o ambiguo en las notas originales (texto literal no verificado)
\end{norma}

\begin{table}[H]
\centering
\small
\caption{Comparación entre los arts.~118 y 119}
\begin{tabularx}{\textwidth}{>{\raggedright\arraybackslash}p{0.17\textwidth} Y Y}
\toprule
\textbf{Aspecto} & \textbf{Art.~118 · de pleno derecho} & \textbf{Art.~119 · por conocimiento} \\
\midrule
\textbf{Prueba} & No se necesita demostrar nada. & Se requiere que el tercero conociera la cesación de pagos. \\
\textbf{Declaración} & El juez la detecta y declara de oficio. & Se discute en la acción de ineficacia iniciada por el síndico. \\
\textbf{Actos} & A título gratuito; pagos anticipados; garantía real a un quirografario. & Actos que afectan a todos los acreedores (por ejemplo, venta de un inmueble al socio cuyo producido no reingresa al patrimonio). \\
\textbf{Si el dinero reingresa} & --- & No hay afectación de acreedores (se compra otro inmueble o se deposita para pagar): no hay acto ineficaz. \\
\bottomrule
\end{tabularx}
\end{table}

\section{La acción de ineficacia (art.~120)}

\subsection{Legitimación del síndico y mayoría de acreedores}

La acción de ineficacia por conocimiento del estado de cesación de pagos debe ser iniciada por el \textbf{síndico}, dentro de la quiebra. La ley exige la aceptación de la \textbf{mayoría absoluta de los acreedores} para que el síndico inicie la acción.

\subsection{El cambio en la práctica: el silencio como aceptación}

Hasta hace unos años, el juez enviaba una cédula de notificación a cada acreedor para que prestara su conformidad. Los acreedores no necesitan patrocinio letrado para actuar en un proceso concursal y, aun cuando lo tuvieran, en su gran mayoría no se presentaban a dar su aceptación; por eso las acciones eran muy difíciles de iniciar.

Los jueces advirtieron esta situación y cambiaron el sentido de la notificación: se informa a los acreedores que el síndico pretende iniciar una acción de ineficacia y que, si algún acreedor no acepta que lo haga, debe presentarse. El \textbf{silencio} se considera entonces manifestación de voluntad en cuanto a la aceptación, y hoy solo se presentan los acreedores que no quieren que la acción se inicie. Con este cambio comenzaron a funcionar las acciones de ineficacia.

\subsection{Legitimación subsidiaria de los acreedores}

El art.~120 habilita además a los acreedores que quieran iniciar la acción cuando el síndico no lo hace: pueden pedir que se intime al síndico y, si este no actúa, promoverla ellos mismos.

\begin{norma}{Art.~120 · Acción por los acreedores}
Contenido expuesto: si el síndico no inicia la acción de ineficacia, los acreedores interesados pueden pedir que se lo intime a hacerlo y, si no lo hace, promoverla ellos mismos.

\par\smallskip\textit{Texto literal pendiente de verificación: se resume lo expuesto; cotejar con la versión vigente de la Ley 24.522.}
% TODO: contenido incompleto o ambiguo en las notas originales (texto literal no verificado)
\end{norma}

\begin{important}
\textbf{Período de sospecha:} desde el inicio de la cesación de pagos hasta el decreto de quiebra. \textbf{Art.~118:} ineficacia de pleno derecho (el juez actúa de oficio). \textbf{Art.~119:} ineficacia por conocimiento del tercero (acción del síndico, con la mayoría de los acreedores). \textbf{Art.~120:} habilita a los acreedores si el síndico no actúa.
\end{important}

\subsection*{Cuadro Cornell del capítulo}

\begin{longtable}{>{\raggedright\arraybackslash}p{0.27\textwidth} >{\raggedright\arraybackslash}p{0.64\textwidth}}
\toprule
\textbf{Preguntas / palabras clave} & \textbf{Notas / respuestas} \\
\midrule
\endfirsthead
\toprule
\textbf{Preguntas / palabras clave} & \textbf{Notas / respuestas} \\
\midrule
\endhead
\textbf{¿Qué es el período de sospecha?} & El lapso que va desde el inicio de la cesación de pagos hasta el decreto de quiebra, durante el cual los actos del deudor pueden considerarse ineficaces. \\[0.6em]
\textbf{¿Qué actos son ineficaces de pleno derecho (art.~118)?} & Los actos a título gratuito, los pagos anticipados de créditos aún no vencidos y la constitución de garantía real a un acreedor quirografario, que pasa a ser privilegiado. No se necesita probar nada y el juez actúa de oficio. \\[0.6em]
\textbf{¿Cuándo importa lo que sabía el tercero (art.~119)?} & Cuando el tercero que intervino en el acto conocía la cesación de pagos del deudor y el acto afecta a todos los acreedores. Si el producido reingresa al patrimonio (por ejemplo, se compra otro inmueble o se deposita para pagar), no hay afectación ni acto ineficaz. \\[0.6em]
\textbf{¿Quién puede iniciar la acción de ineficacia?} & El síndico, con la aceptación de la mayoría absoluta de los acreedores; en la práctica actual, el silencio ante la notificación se considera aceptación. Si el síndico no actúa, los acreedores pueden intimarlo y, luego, promover la acción (art.~120). \\
\midrule
\multicolumn{2}{>{\raggedright\arraybackslash}p{\cornellfull}}{\textbf{Resumen:} los actos del deudor realizados entre el inicio de la cesación de pagos y el decreto de quiebra pueden caer. Son ineficaces de pleno derecho los gratuitos, los pagos anticipados y la garantía real a un quirografario (art.~118); los demás actos perjudiciales lo son si el tercero conocía la cesación de pagos (art.~119). La acción corresponde al síndico y, subsidiariamente, a los acreedores (art.~120).} \\
\bottomrule
\end{longtable}

% ------------------------------------------------------------
\chapter[Relaciones jurídicas preexistentes]{Efectos sobre las relaciones jurídicas preexistentes}
% ------------------------------------------------------------

\section{Efectos generales sobre las relaciones jurídicas}

En cuanto a los efectos generales, todos los acreedores por título o causa anterior a la fecha del decreto de quiebra quedan incluidos en la quiebra. Los \textbf{intereses se suspenden}. Si el deudor tuviera bienes de terceros, deben restituirse; si fuese al revés, los acreedores también pueden plantear una acción.

\begin{norma}{Art.~129 · Suspensión de intereses}
«La declaración de quiebra suspende el curso de intereses de todo tipo. Sin embargo, los compensatorios devengados con posterioridad que correspondan a créditos amparados con garantías reales pueden ser percibidos hasta el límite del producido del bien gravado después de pagadas las costas, los intereses preferidos anteriores a la quiebra y el capital. Asimismo, tampoco se suspenden los intereses compensatorios devengados con posterioridad que correspondan a créditos laborales.»

\par\smallskip\textit{Norma relacionada (se indica que «los intereses se suspenden», sin citar el número). Texto literal verificado. Ley 24.522 según Ley 26.684 (B.O. 30/06/2011).}
\end{norma}

\section{Contratos en curso de ejecución (arts.~143 a 145)}

\subsection{Artículo 143: quién cumplió}

Cuando hay contratos en curso de ejecución, debe establecerse qué parte cumplió. Si el \textbf{fallido cumplió} con su prestación y el cocontratante no, la quiebra (el síndico) le exige al cocontratante que cumpla. Si, al revés, el \textbf{cocontratante cumplió} y el fallido no, el cocontratante debe presentarse a verificar su crédito.

\begin{norma}{Art.~143 · Contratos en curso de ejecución}
Contenido expuesto: (a) si el fallido cumplió y el cocontratante no, la quiebra (el síndico) le exige que cumpla; (b) si el cocontratante cumplió y el fallido no, el cocontratante debe presentarse a verificar su crédito; (c) si hay prestaciones recíprocas pendientes, se remite al art.~144.

\par\smallskip\textit{Texto literal pendiente de verificación: se resume lo expuesto; cotejar con la versión vigente de la Ley 24.522.}
% TODO: contenido incompleto o ambiguo en las notas originales (texto literal no verificado)
\end{norma}

\subsection{Artículo 144: prestaciones recíprocas pendientes}

Cuando el contrato en curso de ejecución tiene \textbf{prestaciones recíprocas pendientes}, corresponde acudir al art.~144 y establecer si en la quiebra se dispone la \textbf{continuación de la explotación de la empresa}. Si se dispone, debe definirse si el contrato se mantiene o no.

\begin{norma}{Art.~144 · Prestaciones recíprocas pendientes}
Contenido expuesto: si hay prestaciones recíprocas pendientes, debe establecerse si se dispone la continuación de la explotación de la empresa; si se continúa, corresponde definir si el contrato se mantiene o no.

\par\smallskip\textit{Texto literal pendiente de verificación: se resume lo expuesto; cotejar con la versión vigente de la Ley 24.522.}
% TODO: contenido incompleto o ambiguo en las notas originales (texto literal no verificado)
\end{norma}

\subsection{Artículo 145: cláusulas de resolución por quiebra}

El art.~145 es considerado un artículo \textbf{discutido o discutible}. Si en el contrato se hubiera establecido que ante la quiebra de una de las partes (o por incumplimiento) el contrato queda resuelto, esa \textbf{cláusula es inaplicable}. Los contratos con prestaciones personales del fallido (\textit{intuitu personae}) se \textbf{extinguen}.

\begin{norma}{Art.~145 · Cláusulas de resolución por quiebra}
Contenido expuesto: la cláusula contractual que prevé la resolución del contrato por la quiebra de una de las partes (o por incumplimiento) es inaplicable. Los contratos con prestaciones personales (\textit{intuitu personae}) del fallido se extinguen.

\par\smallskip\textit{Texto literal pendiente de verificación: se resume lo expuesto; cotejar con la versión vigente de la Ley 24.522.}
% TODO: contenido incompleto o ambiguo en las notas originales (texto literal no verificado; se indica que el artículo es discutible)
\end{norma}

\begin{note}
Frente a la consulta sobre si podría pactarse que el contrato \emph{no} se resuelva por la quiebra, se explicó que lo habitual es lo inverso: quien contrata suele prever que, ante el incumplimiento o la quiebra de alguna de las partes, el contrato queda resuelto, pues no se querrá mantener un contrato con quien está en quiebra y no va a pagar. Esa cláusula no tiene efectos en la quiebra, ya que interviene el juez, que resolverá el contrato si carece de sentido. Se invitó a pensar el caso desde el lado del cocontratante cuando es el fallido quien debe cumplir y no cumple.
% TODO: contenido incompleto o ambiguo en las notas originales (la explicación sobre el alcance de la cláusula no queda del todo precisada)
\end{note}

\begin{table}[H]
\centering
\small
\caption{Contratos en curso de ejecución}
\begin{tabularx}{\textwidth}{>{\raggedright\arraybackslash}p{0.38\textwidth} Y}
\toprule
\textbf{Situación} & \textbf{Consecuencia} \\
\midrule
\textbf{Cumplió el fallido; no cumplió el cocontratante} & El síndico exige al cocontratante que cumpla. \\
\textbf{Cumplió el cocontratante; no cumplió el fallido} & El cocontratante debe presentarse a verificar su crédito. \\
\textbf{Prestaciones recíprocas pendientes} & Se aplica el art.~144: según se continúe o no la explotación, se define si el contrato se mantiene. \\
\textbf{Cláusula de resolución por quiebra (art.~145)} & Inaplicable. \\
\textbf{Contrato \textit{intuitu personae} del fallido} & Se extingue. \\
\bottomrule
\end{tabularx}
\end{table}

\section{Contratos de locación (art.~157)}

El art.~157 contempla los distintos supuestos de locación. Si el fallido es \textbf{locador}, el contrato se mantiene, porque los alquileres ingresan a la quiebra. Si el fallido es \textbf{locatario}, debe distinguirse: si se trata de \textbf{su vivienda}, queda fuera de la quiebra; si es una \textbf{locación comercial} en la que opera comercialmente, debe esperarse a que se determine si se continúa o no la explotación de la empresa.

\begin{norma}{Art.~157 · Locación}
Contenido expuesto: si el fallido es locador, el contrato se mantiene (los alquileres ingresan a la quiebra). Si es locatario de su vivienda, queda fuera de la quiebra; si la locación es comercial, se espera la decisión sobre la continuación de la explotación.

\par\smallskip\textit{Texto literal pendiente de verificación: se resume lo expuesto; cotejar con la versión vigente de la Ley 24.522.}
% TODO: contenido incompleto o ambiguo en las notas originales (texto literal no verificado)
\end{norma}

\begin{table}[H]
\centering
\small
\caption{Locación y quiebra}
\begin{tabularx}{\textwidth}{>{\raggedright\arraybackslash}p{0.25\textwidth} >{\raggedright\arraybackslash}p{0.25\textwidth} Y}
\toprule
\textbf{Posición del fallido} & \textbf{Tipo de locación} & \textbf{Consecuencia} \\
\midrule
\textbf{Locador} & --- & El contrato se mantiene; los alquileres entran en la quiebra. \\
\textbf{Locatario} & De su vivienda & Queda fuera de la quiebra. \\
\textbf{Locatario} & Comercial & Depende de si se dispone la continuación de la explotación. \\
\bottomrule
\end{tabularx}
\end{table}

\section{Efectos sobre el contrato de trabajo}

Los contratos de trabajo se \textbf{suspenden por 60 días corridos}. Esta suspensión depende de lo que se resuelva sobre la continuación de la explotación de la empresa: si la empresa continúa, el contrato se retoma y continúa; si se decide no continuar la explotación, el contrato se extingue.

\begin{important}
\textbf{Contrato de trabajo:} suspensión por 60 días corridos. Si la empresa continúa, el contrato se retoma; si no continúa, se extingue.
\end{important}

\subsection*{Cuadro Cornell del capítulo}

\begin{longtable}{>{\raggedright\arraybackslash}p{0.27\textwidth} >{\raggedright\arraybackslash}p{0.64\textwidth}}
\toprule
\textbf{Preguntas / palabras clave} & \textbf{Notas / respuestas} \\
\midrule
\endfirsthead
\toprule
\textbf{Preguntas / palabras clave} & \textbf{Notas / respuestas} \\
\midrule
\endhead
\textbf{¿Qué ocurre con los acreedores anteriores y los intereses?} & Los acreedores por título o causa anterior al decreto de quiebra quedan incluidos en ella y los intereses se suspenden (art.~129, con las excepciones allí previstas). Los bienes de terceros en poder del deudor se restituyen. \\[0.6em]
\textbf{¿Qué sucede con los contratos en curso (art.~143)?} & Si cumplió el fallido, el síndico exige el cumplimiento al cocontratante; si cumplió el cocontratante, este debe verificar su crédito; con prestaciones recíprocas pendientes se remite al art.~144. \\[0.6em]
\textbf{¿Qué define el art.~144?} & Si se dispone la continuación de la explotación de la empresa y, en ese caso, si el contrato se mantiene o no. \\[0.6em]
\textbf{¿Qué efecto tienen las cláusulas de resolución por quiebra (art.~145)?} & Son inaplicables. Los contratos \textit{intuitu personae} del fallido se extinguen. \\[0.6em]
\textbf{¿Qué pasa con los alquileres del fallido (art.~157)?} & Si es locador, el contrato se mantiene y los alquileres entran en la quiebra. Si es locatario de su vivienda, queda fuera de la quiebra; si es comercial, se espera la decisión sobre la continuación. \\[0.6em]
\textbf{¿Qué pasa con el contrato de trabajo?} & Se suspende por 60 días corridos; si la empresa continúa se retoma y si no continúa se extingue. \\
\midrule
\multicolumn{2}{>{\raggedright\arraybackslash}p{\cornellfull}}{\textbf{Resumen:} los acreedores anteriores quedan incluidos en la quiebra y los intereses se suspenden. Los contratos en curso siguen reglas según quién cumplió (art.~143), la continuación de la empresa (art.~144) y la inaplicabilidad de las cláusulas resolutorias por quiebra (art.~145). La locación depende de la posición del fallido (art.~157) y el contrato de trabajo se suspende por 60 días.} \\
\bottomrule
\end{longtable}

% ------------------------------------------------------------
\chapter[Continuación de la explotación]{Continuación de la explotación de la empresa}
% ------------------------------------------------------------

\section{Concepto y requisitos (arts.~189 y siguientes)}

\subsection{Un instituto central: se conserva la empresa, no el empresario}

La continuación de la explotación de la empresa es un instituto muy importante, muy trabajado por la doctrina y modificado reiteradamente por la legislación. Frente a la pregunta de por qué continuaría una empresa cuya quiebra fue decretada, la respuesta es que \textbf{lo que se conserva es la empresa, no el empresario}: lo que se quiere mantener es la \textbf{actividad}.

\subsection{Socialmente útil y económicamente viable}

La ley contempla la posibilidad de continuar la explotación cuando la empresa es \textbf{socialmente útil} y \textbf{económicamente viable}. Ambos requisitos deben darse para que el juez resuelva la continuidad.

\begin{definition}{Socialmente útil}
La empresa es una fuente de trabajo: lo que se busca es mantener la fuente de trabajo.
\end{definition}

\begin{definition}{Económicamente viable}
La empresa no genera nuevos pasivos, o más pasivos de los que ya tiene. No significa que no haya ningún pasivo nuevo, sino que los que surjan sean propios de la actividad ordinaria, sin que se desborden con pasivos que podrían haberse evitado.
\end{definition}

Si bien en términos de trabajo toda empresa que da empleo es socialmente útil, \textbf{deben darse los dos requisitos}. Mantener la fuente de trabajo sin que los trabajadores cobren su sueldo tampoco es la respuesta que se busca. Por eso se realizan numerosos análisis previos para decidir si se continúa o no con la explotación.

\begin{important}
Se conserva la empresa (la actividad), no el empresario. Requisitos acumulativos: que sea \textbf{socialmente útil} (fuente de trabajo) y \textbf{económicamente viable} (que no genere nuevos pasivos, o que sean solo los propios de la actividad ordinaria).
\end{important}

\section{La regla actual y la venta de la empresa como unidad (art.~204)}

\subsection{Hoy la continuación es la regla}

Actualmente la continuación de la explotación es la \textbf{regla} y no la excepción, como ocurría antes; este criterio fue modificado. Ejemplo de ello son las cooperativas de trabajo que se quedaron con las empresas, lo que antes se denominaba \textbf{empresas recuperadas}.

\subsection{Prelación en la liquidación: vender la empresa como unidad}

El art.~204 establece la prioridad o prelación en la liquidación de los bienes, dado que la quiebra es un proceso de liquidación. Lo primero que debe intentarse es \textbf{vender la empresa como unidad}: una empresa que está en marcha permite obtener mejor resultado. En este punto aparecen nuevamente los trabajadores, que en su gran mayoría son quienes se quedan con esas empresas.

\begin{norma}{Art.~204 · Prelación en la enajenación}
Contenido expuesto: en la liquidación, la primera opción es vender la empresa como unidad, en marcha, porque se obtiene mejor resultado; luego se prevén otras alternativas, que se estudiarán posteriormente.

\par\smallskip\textit{Texto literal pendiente de verificación: se resume lo expuesto; cotejar con la versión vigente de la Ley 24.522.}
% TODO: contenido incompleto o ambiguo en las notas originales (texto literal no verificado; alternativas posteriores no desarrolladas)
\end{norma}

\begin{example}{una cadena de restaurantes convertida en cooperativas}
La cadena de restaurantes conocida como Los Chanchitos aparece hoy como cooperativas de trabajo, porque los trabajadores se quedaron con ella; son varios locales (tres o cuatro). Ilustra la distinción entre empresa y empresario: se mantiene la actividad en desarrollo y, con ella, las fuentes de trabajo.
\end{example}

\section{Continuación inmediata y continuación mediata}

Existen dos formas de disponer la continuación: la \textbf{inmediata} y la \textbf{mediata}.

\subsection{Continuación inmediata}

La continuación inmediata puede decidirla el \textbf{síndico} o la \textbf{cooperativa de trabajo}. En este último caso se requiere que se unan, por lo menos, \textbf{las dos terceras partes} de los trabajadores; muchas veces la cooperativa aún no está formada y puede presentarse y pedir la continuación como \textbf{cooperativa en formación}.

Para determinar si procede, se considera si la interrupción de la actividad o de la producción podría causar un daño mucho más grave, o si la continuación permite mantener el patrimonio que luego se liquidará. Por ejemplo, si un proceso productivo está casi por terminar y se lo corta a la mitad, lo realizado se pierde; también apagar las máquinas y volver a encenderlas tiene un costo. En estos casos lo más importante es que la continuación se dé de forma inmediata, y se debe informar al juez dentro del plazo de \textbf{24 horas}.

\begin{norma}{Art.~189 (primer párrafo) · Continuación inmediata}
«Continuación inmediata. El síndico puede continuar de inmediato con la explotación de la empresa o alguno de sus establecimientos, si de la interrupción pudiera resultar con evidencia un daño grave al interés de los acreedores y a la conservación del patrimonio, si se interrumpiera un ciclo de producción que puede concluirse o entiende que el emprendimiento resulta económicamente viable. También la conservación de la fuente de trabajo habilita la continuación inmediata de la explotación de la empresa o de alguno de sus establecimientos, si las dos terceras partes del personal en actividad o de los acreedores laborales, organizados en cooperativa, incluso en formación, la soliciten al síndico o al juez, si aquél todavía no se hubiera hecho cargo, a partir de la sentencia de quiebra y hasta cinco (5) días luego de la última publicación de edictos en el diario oficial que corresponda a la jurisdicción del establecimiento. El síndico debe ponerlo en conocimiento del juez dentro de las veinticuatro (24) horas. El juez puede adoptar las medidas que estime pertinentes, incluso la cesación de la explotación, con reserva de lo expuesto en los párrafos siguientes. Para el caso que la solicitud a que refiere el segundo párrafo el presente, sea una cooperativa en formación, la misma deberá regularizar su situación en un plazo de cuarenta (40) días, plazo que podría extenderse si existiesen razones acreditadas de origen ajeno a su esfera de responsabilidad que impidan tal cometido.»

\par\smallskip\textit{Primer párrafo, según art.~16 de la Ley 26.684. Texto literal verificado. Ley 24.522 según Ley 26.684 (B.O. 30/06/2011).}
\end{norma}

\subsection{Continuación mediata}

La continuación mediata es, quizás, la más común. En ella el \textbf{síndico presenta un informe} dentro del plazo de \textbf{20 días corridos} desde que aceptó el cargo, y el juez resuelve.

\begin{table}[H]
\centering
\small
\caption{Continuación inmediata y mediata}
\begin{tabularx}{\textwidth}{>{\raggedright\arraybackslash}p{0.16\textwidth} Y Y}
\toprule
\textbf{Aspecto} & \textbf{Inmediata} & \textbf{Mediata} \\
\midrule
\textbf{Quién la promueve} & El síndico o la cooperativa de trabajo (dos terceras partes de los trabajadores; puede estar en formación). & El síndico, mediante un informe al juez. \\
\textbf{Plazo} & Se informa al juez dentro de las 24 horas. & Informe dentro de los 20 días corridos desde que aceptó el cargo. \\
\textbf{Criterio} & Evitar un daño mayor, no cortar un ciclo de producción, mantener el patrimonio a liquidar. & El juez resuelve a partir del informe del síndico. \\
\textbf{Frecuencia} & No es la más común. & Es la más común. \\
\bottomrule
\end{tabularx}
\end{table}

\section{Continuación en los servicios públicos}

Si quien está en quiebra es una empresa que presta un \textbf{servicio público}, la ley dispone que \textbf{no puede cortarse}, porque se afectaría a todos los usuarios (por ejemplo, la luz o el gas). En esos casos la continuación es inmediata y automática, y prevalece sobre los propios intereses de los acreedores, porque el servicio público es para todos. Debe notificarse a la autoridad que concedió el servicio o a su ente regulador. Aparece entonces la parte administrativa de determinar quién reemplazará a la empresa, que no seguirá prestando el servicio; lo que no puede hacerse es cortarlo.

\section{Informe del síndico y plan de explotación}

\subsection{Contenido del informe}

El síndico debe presentar un informe para recomendar la continuación de la explotación, que incluye un \textbf{plan de explotación}: debe indicar si la explotación es económicamente viable y si puede evitarse generar nuevos pasivos. Esto no significa que no haya ningún pasivo nuevo, sino que sean propios de la actividad ordinaria.

\subsection{Trabajadores y contratos que se mantienen}

El plan debe determinar qué \textbf{trabajadores} se mantendrán (siempre se procura mantener a todos) y cuáles son los \textbf{contratos necesarios} para mantener la empresa. Por ejemplo, en una empresa de impresión con fotocopiadoras alquiladas, si se perdiera ese contrato no podría continuar el negocio; por lo tanto, ese contrato debe mantenerse, y así debe indicarlo el síndico en su informe.

\begin{norma}{Art.~190 · Trámite común para todos los procesos}
«En toda quiebra, aun las comprendidas en el artículo precedente, el síndico debe informar al juez dentro de los veinte (20) días corridos contados a partir de la aceptación del cargo, sobre la posibilidad de continuar con la explotación de la empresa del fallido o de alguno de sus establecimientos y la conveniencia de enajenarlos en marcha. En la continuidad de la empresa se tomará en consideración el pedido formal de los trabajadores en relación de dependencia que representen las dos terceras partes del personal en actividad o de los acreedores laborales quienes deberán actuar en el período de continuidad bajo la forma de una cooperativa de trabajo. A tales fines deberá presentar en el plazo de veinte (20) días, a partir del pedido formal, un proyecto de explotación conteniendo las proyecciones referentes a la actividad económica que desarrollará, del que se dará traslado al síndico para que en plazo de cinco (5) días emita opinión al respecto. El término de la continuidad de la empresa, cualquiera sea su causa, no hace nacer el derecho a nuevas indemnizaciones laborales. El informe del síndico debe expedirse concretamente sobre los siguientes aspectos: 1) La posibilidad de mantener la explotación sin contraer nuevos pasivos, salvo los mínimos necesarios para el giro de la explotación de la empresa o establecimiento; 2) La ventaja que resultaría para los acreedores de la enajenación de la empresa en marcha; 3) La ventaja que pudiere resultar para terceros del mantenimiento de la actividad; 4) El plan de explotación acompañado de un presupuesto de recursos, debidamente fundado; 5) Los contratos en curso de ejecución que deben mantenerse; 6) En su caso, las reorganizaciones o modificaciones que deben realizarse en la empresa para hacer económicamente viable su explotación; 7) Los colaboradores que necesitará para la administración de la explotación; 8) Explicar el modo en que se pretende cancelar el pasivo preexistente. En caso de disidencias o duda respecto de la continuación de la explotación por parte de los trabajadores, el juez, si lo estima necesario, puede convocar a una audiencia a los intervinientes en la articulación y al síndico, para que comparezcan a ella, con toda la prueba de que intenten valerse. El juez, a los efectos del presente artículo y en el marco de las facultades del artículo 274, podrá de manera fundada extender los plazos que se prevén en la ley para la continuidad de la empresa, en la medida que ello fuere razonable para garantizar la liquidación de cada establecimiento como unidad de negocio y con la explotación en marcha.»

\par\smallskip\textit{Según art.~17 de la Ley 26.684. Texto literal verificado. Ley 24.522 según Ley 26.684 (B.O. 30/06/2011).}
\end{norma}

\begin{norma}{Art.~191 · Autorización del juez}
«La autorización para continuar con la actividad de la empresa del fallido o de alguno de sus establecimientos será dada por el juez en caso de que de su interrupción pudiera emanar una grave disminución del valor de realización, se interrumpiera un ciclo de producción que puede concluirse, en aquellos casos que lo estime viable económicamente o en resguardo de la conservación de la fuente laboral de los trabajadores de la empresa declarada en quiebra. En su autorización el juez debe pronunciarse explícitamente por lo menos sobre: 1) El plan de la explotación, para lo cual podrá hacerse asesorar por expertos o entidades especializadas; 2) El plazo por el que continuará la explotación; a estos fines se tomará en cuenta el ciclo y el tiempo necesario para la enajenación de la empresa; este plazo podrá ser prorrogado por una sola vez, por resolución fundada; 3) La cantidad y calificación profesional del personal que continuará afectado a la explotación; 4) Los bienes que pueden emplearse; 5) La designación o no de uno o más coadministradores; y la autorización al síndico para contratar colaboradores de la administración; 6) Los contratos en curso de ejecución que se mantendrán; los demás quedarán resueltos; 7) El tipo y periodicidad de la información que deberá suministrar el síndico y, en su caso, el coadministrador o la cooperativa de trabajo. Esta resolución deberá ser dictada dentro de los diez (10) días posteriores a la presentación del informe de la sindicatura previsto en el artículo 190. La resolución que rechace la continuación de la explotación es apelable por el síndico y la cooperativa de trabajo.»

\par\smallskip\textit{Norma relacionada (resolución del juez sobre el informe). Según art.~18 de la Ley 26.684. Texto literal verificado. Ley 24.522 según Ley 26.684 (B.O. 30/06/2011).}
\end{norma}

\section{Cooperativas de trabajo y acreedores con garantía real}

\subsection{Artículo 191 bis: asistencia técnica del Estado}

Cuando son los trabajadores quienes mantienen la empresa, intervienen las \textbf{cooperativas de trabajo}. El \textbf{art.~191 bis}, incorporado por la Ley 26.684 (la última reforma importante de la ley), establece el derecho de los trabajadores a recibir \textbf{asistencia técnica por parte del Estado}.

\begin{norma}{Art.~191 bis · Asistencia técnica a las cooperativas}
«En toda quiebra que se haya dispuesto la continuidad de la explotación de la empresa o de alguno de sus establecimientos por parte de las dos terceras partes del personal en actividad o de los acreedores laborales, organizados en cooperativas, incluso en formación, el Estado deberá brindarle la asistencia técnica necesaria para seguir adelante con el giro de los negocios.»

\par\smallskip\textit{Incorporado por art.~19 de la Ley 26.684. Texto literal verificado. Ley 24.522 según Ley 26.684 (B.O. 30/06/2011).}
\end{norma}

\subsection{Acreedores con garantías reales}

Los acreedores con garantías reales no pueden ejecutar sus bienes si se continúa la explotación de la empresa, hasta tanto exista una conformidad posterior para el remate.

\begin{norma}{Art.~195 · Hipoteca y prenda en la continuación de empresa}
«En caso de continuación de la empresa, los acreedores hipotecarios o prendarios no pueden utilizar el derecho a que se refieren los artículos 126, segunda parte, y 209, sobre los bienes necesarios para la explotación, en los siguientes casos: 1) Cuando los créditos no se hallen vencidos a la fecha de la declaración y el síndico satisfaga las obligaciones posteriores en tiempo debido; 2) Cuando los créditos se hallen vencidos a la fecha de la declaración, mientras no cuenten con resolución firme que acredite su calidad de acreedor hipotecario o prendario; 3) Cuando exista conformidad del acreedor hipotecario o prendario para la suspensión de la ejecución. Son nulos los pactos contrarios a las disposiciones de los incisos 1) y 2). Por decisión fundada y a pedido de la cooperativa de trabajadores, el juez de la quiebra podrá suspender las ejecuciones hipotecarias y/o prendarias por un plazo de hasta dos (2) años.»

\par\smallskip\textit{Norma relacionada (no se cita el número). Según art.~21 de la Ley 26.684. Texto literal verificado. Ley 24.522 según Ley 26.684 (B.O. 30/06/2011).}
\end{norma}

\subsection*{Cuadro Cornell del capítulo}

\begin{longtable}{>{\raggedright\arraybackslash}p{0.27\textwidth} >{\raggedright\arraybackslash}p{0.64\textwidth}}
\toprule
\textbf{Preguntas / palabras clave} & \textbf{Notas / respuestas} \\
\midrule
\endfirsthead
\toprule
\textbf{Preguntas / palabras clave} & \textbf{Notas / respuestas} \\
\midrule
\endhead
\textbf{¿Por qué puede continuar una empresa que quebró?} & Porque lo que se conserva es la empresa (la actividad) y no el empresario. Hoy la continuación es la regla y no la excepción. \\[0.6em]
\textbf{¿Qué requisitos exige la continuación?} & Que la empresa sea socialmente útil (fuente de trabajo) y económicamente viable (no genere nuevos pasivos, o solo los propios de la actividad ordinaria). Deben darse ambos. \\[0.6em]
\textbf{¿Qué establece el art.~204?} & La prelación en la liquidación: lo primero es intentar vender la empresa como unidad, en marcha, porque se obtiene mejor resultado. \\[0.6em]
\textbf{¿Cómo se diferencian la continuación inmediata y la mediata?} & Inmediata: la decide el síndico o la cooperativa de trabajo (dos terceras partes del personal, incluso en formación) y se informa al juez dentro de 24 horas. Mediata: el síndico presenta un informe en 20 días corridos desde que aceptó el cargo y el juez resuelve (arts.~189 a 191). \\[0.6em]
\textbf{¿Qué ocurre si quiebra un servicio público?} & No puede cortarse: la continuación es inmediata y automática, prevalece sobre el interés de los acreedores y se notifica a la autoridad concedente o al ente regulador. \\[0.6em]
\textbf{¿Qué debe contener el informe del síndico?} & Un plan de explotación que indique si es viable sin generar pasivos ajenos a la actividad ordinaria, qué trabajadores se mantienen y qué contratos son necesarios (art.~190). \\[0.6em]
\textbf{¿Qué apoyo reciben las cooperativas de trabajo y qué pasa con los acreedores con garantía real?} & Asistencia técnica del Estado (art.~191 bis, Ley 26.684). Los acreedores con garantía real no pueden ejecutar los bienes necesarios para la explotación mientras continúa, salvo las excepciones legales (art.~195). \\
\midrule
\multicolumn{2}{>{\raggedright\arraybackslash}p{\cornellfull}}{\textbf{Resumen:} se conserva la empresa, no el empresario. La continuación procede si la actividad es socialmente útil y económicamente viable y hoy es la regla; la primera opción de liquidación es vender la empresa en marcha (art.~204). Puede disponerse de modo inmediato (síndico o cooperativa, aviso al juez en 24 horas) o mediato (informe del síndico en 20 días), con plan de explotación; en los servicios públicos es automática. Las cooperativas de trabajo reciben asistencia técnica del Estado y los acreedores con garantía real no pueden ejecutar los bienes necesarios mientras la explotación continúa.} \\
\bottomrule
\end{longtable}

% ------------------------------------------------------------
\chapter[Extensión de quiebra]{Extensión de quiebra}
% ------------------------------------------------------------

\section{Concepto de extensión de quiebra}

La extensión de quiebra constituye una \textbf{excepción al presupuesto objetivo de cesación de pagos}: hay personas a las que se les decreta la quiebra sin que estén en cesación de pagos. La ley prevé los casos en dos artículos importantes: el \textbf{art.~160} y el \textbf{art.~161}.

Se distingue entre la \textbf{quiebra refleja} (art.~160) y la \textbf{extensión de quiebra propiamente dicha} (art.~161).

\section{Quiebra refleja (art.~160)}

La quiebra refleja es aquella por la cual, decretada la quiebra de un sujeto, se la extiende a los \textbf{socios con responsabilidad ilimitada}. Hoy, en la práctica, no se ven muchas sociedades colectivas, ni tantas sociedades en comandita (que tienen dos clases de socios), por lo que el art.~160 \textbf{quedó muy limitado}. Si hubiera socios con responsabilidad ilimitada, a la sociedad se le decreta la quiebra y también a los socios.

\begin{norma}{Art.~160 · Socios con responsabilidad ilimitada}
«La quiebra de la sociedad importa la quiebra de sus socios con responsabilidad ilimitada. También implica la de los socios con igual responsabilidad que se hubiesen retirado o hubieren sido excluidos después de producida la cesación de pagos, por las deudas existentes a la fecha en la que el retiro fuera inscripto en el Registro Público de Comercio, justificadas en el concurso.»

\par\smallskip\textit{Texto según Ley 24.522 (texto ordenado, Infoleg).}
\end{norma}

\section{Extensión de quiebra propiamente dicha (art.~161)}

El art.~161 prevé \textbf{tres supuestos}.

\subsection{Primer supuesto: actuación en interés personal}

El inciso 1 se refiere a un sujeto que realiza actos como si los realizara a nombre de la persona fallida, pero en realidad lo hace en su \textbf{propio interés}. En ese caso, la ley dispone que ese sujeto también debe estar en quiebra.

\begin{example}{actuación en interés personal}
Hay mercadería dentro de la actividad; una persona la saca y la vende como si fuese la sociedad quien vende, pero en realidad lo hace por su propio interés: se guarda el dinero en el bolsillo. Esa conducta puede conducir a que esa persona sea declarada en quiebra.
\end{example}

\subsection{Segundo supuesto: abuso de control}

El inciso 2 trata el \textbf{abuso de control}. Cuando hay una sociedad controlante y otra controlada, y la controlante, a partir de decisiones que pudo adoptar por su posición de control, terminó produciendo la quiebra de la controlada, la ley dispone que la controlante también esté en quiebra.

\subsection{Tercer supuesto: confusión patrimonial}

El inciso 3 se refiere a la \textbf{confusión patrimonial}: cuando hay un patrimonio inescindible, que no puede separarse, y no pueden identificarse los activos y pasivos que corresponden a cada uno de varios sujetos.

\subsection{Finalidad de la extensión}

La extensión se decreta para \textbf{tener más activos para liquidar} y poder pagar a los acreedores de la quiebra principal: se amplía la quiebra a los bienes o al patrimonio de las sociedades alcanzadas.

\begin{norma}{Art.~161 · Actuación en interés personal. Controlantes. Confusión patrimonial}
«La quiebra se extiende: 1) A toda persona que, bajo la apariencia de la actuación de la fallida, ha efectuado los actos en su interés personal y dispuesto de los bienes como si fueran propios, en fraude a sus acreedores; 2) A toda persona controlante de la sociedad fallida, cuando ha desviado indebidamente el interés social de la controlada, sometiéndola a una dirección unificada en interés de la controlante o del grupo económico del que forma parte; 3) A toda persona respecto de la cual existe confusión patrimonial inescindible, que impida la clara delimitación de sus activos y pasivos o de la mayor parte de ellos.»

\par\smallskip\textit{Texto según Ley 24.522 (texto ordenado, Infoleg). Se omiten los párrafos finales del artículo.}
\end{norma}

\begin{table}[H]
\centering
\small
\caption{Quiebra refleja y extensión propiamente dicha (arts.~160 y 161)}
\begin{tabularx}{\textwidth}{>{\raggedright\arraybackslash}p{0.19\textwidth} >{\raggedright\arraybackslash}p{0.10\textwidth} >{\raggedright\arraybackslash}p{0.23\textwidth} Y}
\toprule
\textbf{Supuesto} & \textbf{Art.} & \textbf{Quién queda alcanzado} & \textbf{Ejemplo o precisión} \\
\midrule
\textbf{Quiebra refleja} & 160 & Socios con responsabilidad ilimitada & Hoy casi no hay sociedades colectivas; el art.~160 «quedó muy limitado». \\
\textbf{Actuación en interés personal} & 161 inc.~1 & Quien actúa a nombre del fallido pero en su propio interés & Vende mercadería como si fuera la sociedad y se guarda el dinero. \\
\textbf{Abuso de control} & 161 inc.~2 & La sociedad controlante & La controlante, por su posición de control, produce la quiebra de la controlada. \\
\textbf{Confusión patrimonial} & 161 inc.~3 & Los sujetos con patrimonio inescindible & No se pueden identificar los activos y pasivos de cada uno. \\
\bottomrule
\end{tabularx}
\end{table}

\section{Competencia y verificación de créditos en la quiebra}

\subsection{Juez competente}

La extensión de quiebra la decreta \textbf{el mismo juez que decretó la quiebra principal}. Una vez decretadas las extensiones, el juez que interviene en todas las quiebras (principal y extensiones) es aquel donde se encuentra el sujeto con \textbf{mayor activo}.

\subsection{Períodos informativos y verificación}

En la quiebra se presentan los distintos períodos informativos: verificación de crédito, informe individual e informe general. Debe distinguirse si los acreedores provienen del concurso preventivo, en el caso de una \textbf{quiebra indirecta}:

\begin{itemize}
    \item Si el acreedor verificó su crédito en el concurso y no tiene más deuda que la verificada, \textbf{no debe presentarse a verificar en la quiebra}: el mismo juez actualiza los intereses y el acreedor ya quedó verificado.
    \item Si verificó en el concurso pero tiene \textbf{deuda postconcursal} (deuda posterior a la presentación del concurso, que no correspondía verificar en ese momento), debe presentarse en la quiebra a verificar ese crédito, aclarando que tenía un crédito anterior.
    \item Si solo tiene deuda postconcursal, se presenta a verificar en la quiebra.
\end{itemize}
% TODO: contenido incompleto o ambiguo en las notas originales (en el audio se oye «solamente deuda concursal»; el apunte lo corrige a «postconcursal» por el contexto)

\subsection{Dos modalidades de verificación}

El juez puede determinar dos tipos de verificación: la verificación del \textbf{art.~32}, en la que el acreedor se presenta ante el síndico, o la verificación mediante un \textbf{incidente de verificación}. No se trata de verificación tardía, sino solamente de incidente de verificación.

\begin{norma}{Art.~32 · Verificación de créditos (mencionado)}
Citado como la modalidad de verificación en la que el acreedor se presenta ante el síndico, a diferencia del incidente de verificación que puede disponer el juez en la quiebra.

\par\smallskip\textit{Texto literal pendiente de verificación: se resume lo expuesto; cotejar con la versión vigente de la Ley 24.522.}
% TODO: contenido incompleto o ambiguo en las notas originales (art. 32 solo mencionado)
\end{norma}

\subsection{Verificación tardía en la quiebra}

La verificación tardía es posible dentro de la quiebra hasta la presentación del \textbf{informe de distribución final}, que se estudiará en clases posteriores. Si se la hace antes, el acreedor ingresa como un acreedor más. Si se la hace después de ese informe final, también es posible, pero el acreedor solo cobrará si hay resultado positivo sobre los \textbf{dividendos futuros de distribuciones complementarias}, es decir, si hubieran ingresado nuevos bienes.

\section{Créditos de los trabajadores y pronto pago}

Frente a la consulta sobre qué ocurre con los créditos de los trabajadores en la verificación tardía, o si ya están contemplados previamente, se explicó lo siguiente. En relación con la quiebra, habrá que ver primero si hay continuidad de la explotación de la empresa: si el contrato continúa, los trabajadores deberían cobrar, y ese crédito sería postconcursal, por lo que no entraría en la verificación. Si la quiebra no continuó, tampoco habrá crédito postconcursal.
% TODO: contenido incompleto o ambiguo en las notas originales (pasaje con fragmento inaudible)

En todos los casos, los trabajadores tienen el \textbf{privilegio de pronto pago} dentro de la quiebra, por lo que \textbf{no hace falta que verifiquen su crédito}: pueden ir por el pronto pago. Se señaló que existe un artículo sobre el pronto pago en la quiebra, que puede ser más complejo, pero los trabajadores nunca dejan de tener privilegio. Aun con pronto pago y privilegio, no necesariamente cobran el 100~\% de su crédito.
% TODO: contenido incompleto o ambiguo en las notas originales (no se cita el artículo del pronto pago)

\subsection*{Cuadro Cornell del capítulo}

\begin{longtable}{>{\raggedright\arraybackslash}p{0.27\textwidth} >{\raggedright\arraybackslash}p{0.64\textwidth}}
\toprule
\textbf{Preguntas / palabras clave} & \textbf{Notas / respuestas} \\
\midrule
\endfirsthead
\toprule
\textbf{Preguntas / palabras clave} & \textbf{Notas / respuestas} \\
\midrule
\endhead
\textbf{¿Qué es la extensión de quiebra y para qué sirve?} & Es una excepción al presupuesto objetivo de cesación de pagos: se decreta la quiebra a sujetos que pueden no estar en cesación. Se hace para sumar activos a liquidar y pagar a los acreedores de la quiebra principal. \\[0.6em]
\textbf{¿Qué es la quiebra refleja (art.~160)?} & La extensión de la quiebra de la sociedad a los socios con responsabilidad ilimitada. Hoy quedó muy limitada por la escasez de sociedades colectivas y comanditas. \\[0.6em]
\textbf{¿Qué supuestos prevé el art.~161?} & Actuación en interés personal bajo la apariencia de actuar por el fallido; abuso de control por parte de la controlante; y confusión patrimonial inescindible. \\[0.6em]
\textbf{¿Qué juez interviene en la extensión?} & El que decretó la quiebra principal decreta la extensión; luego interviene en todas las quiebras el juez del sujeto con mayor activo. \\[0.6em]
\textbf{¿Cómo verifican los acreedores en una quiebra indirecta?} & Si ya verificaron en el concurso y no tienen más deuda, no deben verificar de nuevo; la deuda postconcursal debe verificarse en la quiebra. La verificación puede ser por el art.~32 (ante el síndico) o por incidente de verificación. \\[0.6em]
\textbf{¿Qué pasa con la verificación tardía?} & Es posible hasta el informe de distribución final; después de ese informe solo se cobra sobre dividendos futuros de distribuciones complementarias, si ingresan nuevos bienes. \\[0.6em]
\textbf{¿Qué cobran los trabajadores?} & Conservan el privilegio de pronto pago sin necesidad de verificar, aunque no necesariamente cobran el 100~\% de su crédito. \\
\midrule
\multicolumn{2}{>{\raggedright\arraybackslash}p{\cornellfull}}{\textbf{Resumen:} la extensión de quiebra permite alcanzar a socios con responsabilidad ilimitada (art.~160) y a quienes actuaron en interés personal, abusaron del control o confundieron patrimonios (art.~161), con el fin de sumar activos a la quiebra principal. La decreta el juez de la quiebra principal y luego interviene el del mayor activo. En la verificación importa distinguir deuda verificada en el concurso y deuda postconcursal; los trabajadores conservan el pronto pago.} \\
\bottomrule
\end{longtable}

% ------------------------------------------------------------
\chapter[Repaso integrador]{Repaso integrador}
% ------------------------------------------------------------

\begin{longtable}{>{\raggedright\arraybackslash}p{0.27\textwidth} >{\raggedright\arraybackslash}p{0.64\textwidth}}
\toprule
\textbf{Preguntas / palabras clave} & \textbf{Notas / respuestas} \\
\midrule
\endfirsthead
\toprule
\textbf{Preguntas / palabras clave} & \textbf{Notas / respuestas} \\
\midrule
\endhead
\textbf{¿Qué efectos tiene la quiebra decretada sobre el fallido y sus actos?} & El fallido queda inhabilitado y debe cooperar; pierde la administración de sus bienes (desapoderamiento); y sus actos del período de sospecha pueden caerse, sea de pleno derecho (art.~118) o por el conocimiento del tercero (arts.~119 y 120). \\[0.6em]
\textbf{¿Qué pasa con los contratos vigentes y con el trabajo?} & Los contratos siguen reglas propias (arts.~143 a 145 y 157) y el contrato de trabajo se suspende por 60 días. \\[0.6em]
\textbf{¿Cuál es la idea central de la continuación de la explotación?} & Se conserva la empresa y no el empresario: si la actividad es socialmente útil y económicamente viable, puede continuar, de inmediato o previo informe del síndico, a menudo en manos de una cooperativa de trabajo asistida por el Estado. La venta de la empresa en marcha es la primera opción de liquidación. \\[0.6em]
\textbf{¿Para qué sirve la extensión de quiebra?} & Para alcanzar a socios ilimitadamente responsables y a quienes actuaron en interés personal, abusaron del control o confundieron patrimonios, con el fin de sumar activos a la quiebra principal. \\
\midrule
\multicolumn{2}{>{\raggedright\arraybackslash}p{\cornellfull}}{\textbf{Síntesis final:} decretada la quiebra, se busca liquidar el patrimonio del fallido con igualdad entre los acreedores y proteger lo que merece conservarse. Por eso el fallido queda inhabilitado y debe cooperar, pierde la administración de sus bienes y sus actos del período de sospecha pueden caerse (arts.~118 a 120). Los contratos vigentes siguen reglas propias (arts.~143 a 145 y 157) y el contrato de trabajo se suspende por 60 días. Se conserva la empresa y no el empresario: si la actividad es socialmente útil y económicamente viable, puede continuar, y la venta de la empresa en marcha es la primera opción de liquidación. La extensión de quiebra (arts.~160 y 161) permite sumar activos a la quiebra principal. Los trabajadores conservan el pronto pago, aunque no necesariamente cobren todo.} \\
\bottomrule
\end{longtable}

\end{document}
